\documentclass[conference]{IEEEtran}
\usepackage{amsmath,amssymb}
\usepackage{graphicx}
\usepackage{booktabs}
\usepackage{cite}
\usepackage{hyperref}

\begin{document}

\title{Title of the Paper}

\author{
  \IEEEauthorblockN{First Author}
  \IEEEauthorblockA{Department, University\\City, Country\\first@example.org}
  \and
  \IEEEauthorblockN{Second Author}
  \IEEEauthorblockA{Department, University\\City, Country\\second@example.org}
}

\maketitle

\begin{abstract}
State the problem, the approach and the main result in one paragraph.
\end{abstract}

\begin{IEEEkeywords}
first keyword, second keyword, third keyword
\end{IEEEkeywords}

\section{Introduction}
\label{sec:intro}
Introduce the problem and cite the work it builds on \cite{knuth1984texbook, lamport1994latex}. Section~\ref{sec:method}
describes the method and Section~\ref{sec:results} the results.

\section{Method}
\label{sec:method}
Describe the approach. Equations are numbered on the right:
\begin{equation}
  \mathbf{y} = \mathbf{W}\mathbf{x} + \mathbf{b} .
  \label{eq:linear}
\end{equation}

\section{Results}
\label{sec:results}
Table~\ref{tab:main} compares the methods.

\begin{table}[t]
  \centering
  \caption{Main results.}
  \label{tab:main}
  \begin{tabular}{lcc}
    \toprule
    Method & Accuracy & Time (s) \\
    \midrule
    Baseline & 0.81 & 1.0 \\
    Ours & 0.89 & 1.2 \\
    \bottomrule
  \end{tabular}
\end{table}

\section{Conclusion}
Sum up the findings.

\bibliographystyle{IEEEtran}
\bibliography{references}

\end{document}
